% Section 8 (part 3): the full-Stirling series, its rationality, the 1-form and PhiSP.

\subsection{The full-Stirling series, its rationality, the \texorpdfstring{$1$}{1}-form and \texorpdfstring{$\PhiSP$}{the full-Stirling WKB function}}
\label{sec:con-form}

\begin{defn}[The full-Stirling saddle series]
\label{def:gam-c}
Let $\Cstar:=\bbC\setminus\{0,1,\zeta,\zeta^2,\zeta^3,\zeta^4\}$ and $q(L):=(1-L^{-5})/5^5$; no fifth root is taken. For $L\in\Cstar$ fix an open disc $\bbD_L$ around $L$ containing neither $0$ nor any $\zeta^j$, holomorphic branches of $\log(5\sigma)$ and of $\log(\zeta^j-\sigma)$, $j=0,\dots,4$, on $\bbD_L$, and put
\[
\lambda:=-5\log(5L)+\sum_{j=0}^4\log(\zeta^j-L),\qquad\text{so that}\qquad \re^{\lambda}=\frac{\prod_j(\zeta^j-L)}{(5L)^5}=-q(L)
\]
(the letter $\lambda$ is local to this section and \cref{sec:gam}). On $\bbD_L$ define
\begin{align*}
\mathsf S_0(\sigma)&:=\lambda\sigma+5\sigma\log(5\sigma)-5\sigma+\sum_{j=0}^4\bigl[(\zeta^j-\sigma)\log(\zeta^j-\sigma)-(\zeta^j-\sigma)\bigr],\\
\mathsf S_1(\sigma)&:=\tfrac12\log(5\sigma)-\tfrac12\sum_{j=0}^4\log(\zeta^j-\sigma),\qquad
\mathsf S_{2k}(\sigma):=\bk_k\Bigl[(5\sigma)^{1-2k}+\sum_{j=0}^4(\zeta^j-\sigma)^{1-2k}\Bigr]\quad(k\geq1),\\
A_{\mathrm{St}}(\sigma,z)&:=\exp\Bigl(\mathsf S_1(\sigma)+\sum_{k\geq1}z^{2k-1}\mathsf S_{2k}(\sigma)\Bigr).
\end{align*}
Then $\mathsf S_0'(L)=0$ and $\mathsf S_0''=5/\sigma+\sum_j(\zeta^j-\sigma)^{-1}=5/(\sigma(1-\sigma^5))$, which does not vanish at $L$. The \emph{full-Stirling saddle series} is the normalised stationary-phase series (\cref{def:fsp-wick}), with constant term $1$, in the variable $\sigma$ with the measure $\dd\sigma$:
\[
c(z;L):=\SPh_L[\mathsf S_0;A_{\mathrm{St}}](z)=\sum_{n\geq0}c_n(L)z^n,\qquad c_0(L)=1 .
\]
For colour $i$ the series is defined in the same way at the critical point $\sigma=\zeta^iL$, with a disc around $\zeta^iL$, and with the branches and the logarithm of $-q$ chosen there as for $i=0$, so that $\zeta^iL$ is critical.
\end{defn}

With $\log(a/z):=\log a-\log z$ and $\Gst(1+w):=w\,\Gst(w)$ (\cref{def:fsp-stirling}) one has, identically in $\sigma\in\bbD_L$,
\[
\re^{\lambda\sigma/z}\,\Gst(1+5\sigma/z)\prod_{j=0}^4\Gst\bigl((\zeta^j-\sigma)/z\bigr)=(2\pi)^3z^2\exp\Bigl(\mathsf S_0/z+\mathsf S_1+\sum_kz^{2k-1}\mathsf S_{2k}\Bigr),
\]
because the coefficient of $(\log z)/z$ is $-[5\sigma+\sum_j(\zeta^j-\sigma)]=-\sum_j\zeta^j=0$ and that of $\log z$ is $-\frac12+\frac52=2$. So $c(z;L)$ is the normalised saddle series at $\sigma=L$ of the Mellin--Barnes summand $\Gamma(1+5\sigma/z)\prod_j\Gamma((\zeta^j-\sigma)/z)(-q)^{\sigma/z}\,\dd\sigma$, with all six Gamma factors replaced by $\Gst$ and $(-q)^{\sigma/z}:=\re^{\lambda\sigma/z}$; in the Euler integrals the twist factor $\Gamma(1+5\sigma/z)$ carries the measure $\dd y$ and the five factors $\Gamma((\zeta^j-\sigma)/z)$ carry $\dd x/x$.

\emph{Taylor data.} Consider the local power sums $p_k(L):=\sum_j(\zeta^j-L)^{-k}$. Their generating function is
\[
\sum_{k\geq1}p_k(L)\,t^{k-1}=\sum_j(\zeta^j-L-t)^{-1}=\frac{5(L+t)^4}{1-(L+t)^5},
\]
so that $p_k\in(1-L^5)^{-k}\bbQ[L]$, and one has
\begin{gather*}
\mathsf S_0(L+t)-\mathsf S_0(L)=\sum_{k\geq1}\frac{5(-1)^{k+1}L^{-k}+p_k(L)}{k(k+1)}\,t^{k+1},\qquad
-1/\mathsf S_0''(L)=-L(1-L^5)/5,\\
\mathsf S_1(L+t)-\mathsf S_1(L)=\sum_{k\geq1}\frac{(-1)^{k+1}L^{-k}+p_k(L)}{2k}\,t^{k},\\
[t^r]\,\mathsf S_{2k}(L+t)=\bk_k\binom{1-2k}{r}\bigl[5^{1-2k}L^{1-2k-r}+(-1)^rp_{2k-1+r}(L)\bigr].
\end{gather*}

\begin{prop}[Rationality]\genusstatus{AI-proposed, unchecked}
\label{prop:gam-rational}
The series $c(z;L)$ does not depend on the disc $\bbD_L$ nor on the branches of the logarithms, and every $c_n$ is the value at $L$ of an element of $\AL=\bbQ[L,L^{-1},(L^5-1)^{-1}]$. More precisely, $c(z;L)$ depends only on the Taylor coefficients of order at least $2$ of the phase at $L$ and on the normalised amplitude $A_{\mathrm{St}}/A_{\mathrm{St}}(L,0)$; $\lambda$ and the sign convention for $q$ never enter it, they only fix the dictionary $L\leftrightarrow q$.
\end{prop}

\begin{proof}
(1) $\mathsf S_0$ and $A_{\mathrm{St}}$ are holomorphic on $\bbD_L$, $\mathsf S_0'(L)=0$ and $\mathsf S_0''(L)=5/(L(1-L^5))\neq0$. (2) By \cref{lem:fsp-univ}, $c_n(L)$ is a polynomial over $\bbQ$ in the Taylor coefficients at $L$ of $\mathsf S_0$ of orders $3,\dots,2n+2$, in those of $\mathsf S_1-\mathsf S_1(L)$ and of the $\mathsf S_{2k}$ (which determine the normalised amplitude), and in $-1/\mathsf S_0''(L)$. (3) These are the expressions displayed above: polynomials in $L^{-1}$, in the $p_k\in(1-L^5)^{-k}\bbQ[L]$ and in $-L(1-L^5)/5$. (4) Hence $c_n\in\AL$. (5) The Taylor coefficients of order at least $2$ of $\mathsf S_0$, and of order at least $1$ of $\mathsf S_1$ and $\mathsf S_{2k}$, are rational in $\sigma$ and carry no branch, and $\exp(\mathsf S_1(\sigma)-\mathsf S_1(L))$ is the branch equal to $1$ at $L$. Replacing $\log(5\sigma)$ by $\log(5\sigma)+2\pi\ri n$ and $\log(\zeta^j-\sigma)$ by $\log(\zeta^j-\sigma)+2\pi\ri n_j$ changes $\mathsf S_0'$ by $2\pi\ri(5n-\sum_jn_j)$, which the formula for $\lambda$ compensates; then $\mathsf S_0$ changes by the constant $2\pi\ri\sum_j\zeta^jn_j$ and $\mathsf S_1$ by a constant, which the normalised series does not see (\cref{lem:fsp-const}). A logarithm $\lambda+2\pi\ri k$ of $-q$ with $k\neq0$ gives $\mathsf S_0'(L)=2\pi\ri k\neq0$, so $L$ is critical only for $\lambda$; a smaller disc changes nothing. So $c(z;L)$ depends only on $L$, and only through the data named in the statement.
\end{proof}

\begin{lem}[One $1$-form]\genusstatus{AI-proposed, unchecked}
\label{lem:con-form}
Let $\Xi:=-\log(1-u)\in u\bbQ\fps{u}$ and $\cG(x):=\sum_{m\geq1}x^{5m-1}/(m(5m-1))$.
\begin{enumerate}
\item[(a)] In the colour-$0$ summand
\[
q^{\sigma/z}\,\Gamma(1+5\sigma/z)\Big/\prod_{j=0}^4\Gamma\bigl(1+(\sigma-\zeta^j)/z\bigr)
\]
expand $\Gamma(1+5\sigma/z)$ and the five factors $1/\Gamma(1+(\sigma-\zeta^j)/z)$ by Stirling's series, and put $x=1/\sigma$. The exponent is $S(\sigma)/z$ with $S(\sigma):=\sigma\log q+5\sigma\log(5\sigma)-\sum_j(\sigma-\zeta^j)\log(\sigma-\zeta^j)$: the terms in $(\log z)/z$ and the terms linear in $\sigma/z$ cancel. Moreover $S(1/x)=-\cW(x)$ with $\cW(x):=\Xi/x+\cG(x)$, and $\dd\cW/\dd x=-(\Xi+\log(1-x^5))/x^2$; formally, the critical points are $x^5=u$ (an identity of symbols, with $\log(1-x^5)$ and $\log(1-u)$ as series; it is used below only through (c)).
\item[(b)] The half-densities and the measure combine to $(5\sigma/(\sigma^5-1))^{1/2}\,\dd\sigma=-\sqrt5\,(1-x^5)^{-1/2}\,\dd x$. Up to a factor independent of $x$, the summand times $\dd\sigma$ is $\re^{-\cW/z}\fa\,\dd x$ with
\[
\log\fa(x,z)=-\tfrac12\log(1-x^5)+\sum_{k\geq1}\bk_k(xz)^{2k-1}\Bigl[5^{1-2k}-5\sum_{P\geq0}\binom{2k-2+5P}{5P}x^{5P}\Bigr];
\]
$\fa$ does not depend on $u$, $\fa(0,z)=1$, and the coefficient of $x^a$ in $\fa$ is a polynomial in $z$ of degree at most $a$.
\item[(c)] In the hatted variables $x=sX$, $z=s^4\zh$ one has $\cW/z=\hat\cW/\zh$ with
\[
\hat\cW(X;u)=\sum_{m\geq1}\frac{u^{m-1}}m\Bigl[\frac1X+\frac{X^{5m-1}}{5m-1}\Bigr],\qquad \hat\cW'(X)=-\frac1{uX^2}\log\frac{1-uX^5}{1-u};
\]
$X=1$ is critical for every $u$, $\hat\cW''(1)=5/(1-u)$, $\hat\cW(1)=\frac54\tau(u)^{-1}$, and $[t^k]\hat\cW(1+t)=\sum_{m\geq1}u^{m-1}[(-1)^k+\binom{5m-1}{k}/(5m-1)]/m$ lies in $\bbQ\fps{u}$. The hatted amplitude $\hat\fa(X;\zh):=\fa(sX,s^4\zh)$ has
\[
\log\hat\fa=-\tfrac12\log(1-uX^5)+\sum_{k\geq1}\bk_k(uX\zh)^{2k-1}\Bigl[5^{1-2k}-5\sum_{P\geq0}\binom{2k-2+5P}{5P}u^PX^{5P}\Bigr];
\]
it lies in $\bbQ[X,\zh]\fps{u}$, and $\hat\fa(1;0)=(1-u)^{-1/2}$.
\item[(d)] The series $c(z;L)$ of \cref{def:gam-c} (the first Mellin--Barnes summand) and the normalised saddle series at $\sigma=L$ of the second summand of (a), with all Gamma factors replaced by their Stirling forms and the logarithm of $q$ chosen so that $L$ is critical, coincide.
\end{enumerate}
\end{lem}

% Section 8 (part 4): proof of the 1-form lemma; PhiSP and the residue; Theorem P(a), statement and part (i).

\begin{proof}
(a) Replace each $\Gamma(1+y)$ by its Stirling form (\cref{def:fsp-stirling})
\[
\Gst(1+y)=y\,\Gst(y)=\sqrt{2\pi}\,y^{y+1/2}\re^{-y}\exp\Bigl(\sum_{k\geq1}\bk_ky^{1-2k}\Bigr),
\]
with $y=5\sigma/z$ for the numerator and $y=(\sigma-\zeta^j)/z$ for the factors of the denominator, and $\log(a/z)=\log a-\log z$. The terms in $(\log z)/z$ are $-(\log z/z)\,[5\sigma-\sum_j(\sigma-\zeta^j)]=0$, those linear in $\sigma/z$ are $[-5\sigma+\sum_j(\sigma-\zeta^j)]/z=0$, and $q^{\sigma/z}=\re^{\sigma\log q/z}$; what remains of order $1/z$ is $S(\sigma)/z$. At $\sigma=1/x$ write $\log(1/x-\zeta^j)=\log(1-\zeta^jx)-\log x$; with $\sum_j(1/x-\zeta^j)=5/x$ and $\log q+5\log5=\log(1-u)$ this gives $S(1/x)=-\Xi/x-\sum_j(1/x-\zeta^j)\log(1-\zeta^jx)$. Expanding $\log(1-\zeta^jx)=-\sum_n\zeta^{jn}x^n/n$ and summing over $j$, only the terms with $5\mid n$ survive in $\sum_j\log(1-\zeta^jx)$ and only those with $5\mid n+1$ in $\sum_j\zeta^j\log(1-\zeta^jx)$, and
\[
\sum_j(1/x-\zeta^j)\log(1-\zeta^jx)=-\sum_{m\geq1}\frac{x^{5m-1}}m+5\sum_{m\geq1}\frac{x^{5m-1}}{5m-1}=\cG(x).
\]
Then $\cG'(x)=\sum_mx^{5m-2}/m=-x^{-2}\log(1-x^5)$ gives $\dd\cW/\dd x$, which vanishes if and only if $\log(1-x^5)=\log(1-u)$.

(b) The half-densities are $(5\sigma/z)^{1/2}$ from $\Gamma(1+5\sigma/z)$ and $((\sigma-\zeta^j)/z)^{-1/2}$ from the denominator; their product is $z^2(5\sigma/\prod_j(\sigma-\zeta^j))^{1/2}=z^2(5\sigma/(\sigma^5-1))^{1/2}$, and at $\sigma=1/x$ one has $5\sigma/(\sigma^5-1)=5x^4/(1-x^5)$ and $\dd\sigma=-\dd x/x^2$. The Stirling series contribute $\sum_k\bk_k[(5\sigma/z)^{1-2k}-\sum_j((\sigma-\zeta^j)/z)^{1-2k}]=\sum_k\bk_k(xz)^{2k-1}[5^{1-2k}-\sum_j(1-\zeta^jx)^{1-2k}]$, and $\sum_j(1-\zeta^jx)^{-r}=5\sum_{5\mid i}\binom{r-1+i}{i}x^i$ (expand and sum over $j$). The constant $-(2\pi)^{-2}\sqrt5$ and the power $z^2$ do not depend on $x$; what is left is $\re^{-\cW/z}\fa\,\dd x$. Every monomial of $\log\fa$ has $z$-degree at most its $x$-degree, hence so has every monomial of $\fa$.

(c) With $x=sX$ and $z=s^4\zh$ one has $\Xi/(xz)=(\Xi/u)/(X\zh)$ with $\Xi/u=\sum_mu^{m-1}/m$, and $\cG(sX)/(s^4\zh)=\sum_mu^{m-1}X^{5m-1}/(m(5m-1)\zh)$; this is $\hat\cW/\zh$. Termwise in $u$,
\[
\hat\cW'(X)=\sum_m\frac{u^{m-1}}m(X^{5m-2}-X^{-2})=(uX^2)^{-1}\sum_m\frac{u^m}m(X^{5m}-1),
\]
which is the formula stated and vanishes at $X=1$; differentiating the logarithm at $X=1$ gives $\hat\cW''(1)=5/(1-u)$; and $\hat\cW(1)=5\sum_mu^{m-1}/(5m-1)=\frac54\sum_{n\geq0}4u^n/(5n+4)=\frac54\tau^{-1}$ (\cref{def:set-tau}). The Taylor coefficients at $X=1$ follow from those of $(1+t)^{-1}$ and $(1+t)^{5m-1}$. Finally $xz=uX\zh$ and $x^5=uX^5$ in (b).

(d) By \cref{lem:fsp-univ,lem:fsp-const}, a normalised saddle series depends only on the Taylor coefficients of order at least $2$ of the phase at the critical point and on the normalised amplitude. For $y_j:=(\zeta^j-\sigma)/z$ the formal reflection identity \cref{lem:fsp-refl} gives $\log\Gst(y_j)+\log\Gst(1-y_j)=\log2\pi+(y_j-\frac12)(\log y_j-\log(-y_j))$, and $\log y_j-\log(-y_j)=\pi\ri\,\epsilon_j$ with an odd integer $\epsilon_j$ fixed by the branches. Since $1-y_j=1+(\sigma-\zeta^j)/z$, the product $\prod_j\Gst(1+(\sigma-\zeta^j)/z)^{-1}$ is a constant times $\re^{c/z}\re^{\pi\ri(\sum_j\epsilon_j)\sigma/z}\prod_j\Gst((\zeta^j-\sigma)/z)$, with $c$ independent of $\sigma$. Moreover $q^{\sigma/z}=\re^{(\lambda+\pi\ri\epsilon)\sigma/z}$ with an odd integer $\epsilon$, because the logarithms of $q$ and of $-q$ differ by an odd multiple of $\pi\ri$. So the second summand, in Stirling form, is a constant times $\re^{c/z}$ times the first with $\lambda$ replaced by $\lambda+\pi\ri e$, where $e:=\epsilon+\sum_j\epsilon_j$ is even. If $L$ is critical for both, then $e=0$ (step (5) of the proof of \cref{prop:gam-rational}). The two integrands then have the same phase up to an additive constant and the same amplitude up to a constant factor, so their normalised series coincide.
\end{proof}

\begin{defn}[The full-Stirling WKB function and the residue]
\label{def:con-PhiSP}
The \emph{full-Stirling WKB function} is the normalised stationary-phase series, with constant term $1$, of $\re^{-\hat\cW/\zh}\hat\fa\,\dd X$ at $X=1$:
\[
\PhiSP(\zh,u):=\SPh_{X=1}[-\hat\cW;\hat\fa](\zh)=\sum_{n\geq0}\mathrm{ph}^{\mathrm{SP}}_n(u)\zh^n,\qquad\mathrm{ph}^{\mathrm{SP}}_0=1 .
\]
Equivalently,
\[
\PhiSP=\mathrm{SP}_{X=1}\bigl[\re^{-\hat\cW/\zh}\hat\fa\,\dd X\bigr]\Big/\bigl((2\pi\zh/5)^{1/2}\re^{-\hat\cW(1)/\zh}\bigr),
\]
the full stationary phase of \cref{def:fsp-wick} divided by $(2\pi\zh/5)^{1/2}\re^{-\hat\cW(1)/\zh}$, which is the Gaussian factor $(2\pi\zh)^{1/2}\hat\cW''(1)^{-1/2}$ times the critical exponential $\re^{-\hat\cW(1)/\zh}$ times the amplitude value $\hat\fa(1;0)=(1-u)^{-1/2}$; in every comparison below this $\zh$-dependent factor is common to both sides. By \cref{lem:con-form}~(c) and \cref{lem:fsp-univ}, and since the covariance $(1-u)/5$ and the unit $\hat\fa(1;0)=(1-u)^{-1/2}$ lie in $\bbQ\fps u$, every $\mathrm{ph}^{\mathrm{SP}}_n$ lies in $\bbQ\fps{u}$. The \emph{residue} of the flat $1$-form is
\[
\cV:=\Res_{x=0}\re^{-\cW/z}\fa\,\dd x:=\sum_{k\geq1}\frac{(-\Xi/z)^k}{k!}\,[x^{k-1}]\bigl(\re^{-\cG/z}\fa\bigr).
\]
Since $\Xi\in u\bbQ\fps{u}$, $\cG\in x^4\bbQ\fps{x}$ and the coefficient of $x^a$ in $\fa$ has $z$-degree at most $a$, each $u$-order of $\cV$ is a finite sum, and $\cV\in z^{-1}\bbQ[z^{-1}]\fps{u}$.
\end{defn}

\subsection{Theorem P}
\label{sec:con-P}

The ring $\AL$ carries the derivation $\Dq$ with $\Dq L=(L^6-L)/5$; it contains $q=q(L)$, and $\Dq q=q$, so that $\Dq=q\,\dd/\dd q$ on $\AL$. Let $\jmath\colon\AL\to\bbQ\lser{s}$ be the ring homomorphism with $\jmath(L)=s^{-1}$; it is injective, because $\bbQ(L)\to\bbQ\lser{s}$, $L\mapsto s^{-1}$, is a field embedding.

\begin{thm}[Theorem P(a)]\genusstatus{AI-proposed, unchecked}
\label{thm:con-Pa}
Let $c(z;L)=\sum_nc_n(L)z^n$ be the full-Stirling saddle series, with $c_n\in\AL$ (\cref{def:gam-c,prop:gam-rational}).
\begin{enumerate}
\item $\PFfive^{L}\bigl(L\,c(z;L)\bigr)=0$ in $\AL\fps{z}$.
\item In $\bbQ\lser{s}\fps{z}$ one has $\jmath(c)(z)=\PhiSP(s^{-4}z,s^5)$; that is, $\jmath(c_n)=s^{-4n}\mathrm{ph}^{\mathrm{SP}}_n$ for every $n$.
\item $\PhiSP$ solves HPF in $\bbQ\fps{u}\fps{\zh}$.
\item Every $\mathrm{ph}^{\mathrm{SP}}_n$ is a polynomial in $u$; hence $c_n(L)=L^{4n}\mathrm{ph}^{\mathrm{SP}}_n(L^{-5})$ in $\AL$, and $c_n\in\bbQ[L,L^{-1}]$.
\end{enumerate}
\end{thm}

The proof of Theorem~P(a) uses no input from outside the note; part (ii) does not use (iv).

\begin{proof}[Proof of (i)]
In this proof $\sigma=L+t$, $\mathsf S_0$ stands for the formal phase $\mathsf S_0(L+t)-\mathsf S_0(L)\in t^2\AL\fps t$ at $\sigma=L$ (\cref{def:gam-c}; the constant does not matter, \cref{lem:fsp-const}), and $a_0:=A_{\mathrm{St}}/A_{\mathrm{St}}(L,0)=\exp\bigl(\mathsf S_1(L+t)-\mathsf S_1(L)+\sum_kz^{2k-1}\mathsf S_{2k}(L+t)\bigr)\in\AL\fps t\fps z$ for the normalised amplitude, both given by the Taylor data of \cref{def:gam-c}. For $\chi\in\bbQ[\sigma,z]$ put $c\langle\chi\rangle(z;L):=\SPn_L[\mathsf S_0;\chi a_0]$, so that $c=c\langle1\rangle$; the map $\chi\mapsto c\langle\chi\rangle$ extends $\AL[z]$-linearly to $\AL[\sigma,z]$ (\cref{def:fsp-wick}). The letters $\chi$, $a$, $a_0$, $b$, $b_z$, $\mathsf E$, $\mathsf T_z$, $\rho$, $\mathsf l_\rho$, $\mathsf U_\rho$, $\mathsf m_\rho$, $w$, $\Bbbk$ are local to this proof.

\emph{Step 1 (the shift identity).} We apply \cref{lem:fsp-generic}. First, $\mathsf S_0'(L)=0$ and $\mathsf S_0''(\sigma)=5/(\sigma(1-\sigma^5))=-1/\delta(\sigma)$, $\delta(\sigma)=(\sigma^6-\sigma)/5$. Next, $A_{\mathrm{St}}(\sigma,0)^2=\re^{2\mathsf S_1}=5\sigma/\prod_j(\zeta^j-\sigma)=5\sigma/(1-\sigma^5)=-\mathsf r(\sigma)$, so the quotient of $A_{\mathrm{St}}(\sigma,0)$ by $A_{\mathrm{St}}(L,0)$ is $(\mathsf r(\sigma)/\mathsf r(L))^{1/2}$, the root with constant term $1$ at $L$. Finally $b:=\chi\,A_{\mathrm{St}}(\sigma,z)/A_{\mathrm{St}}(\sigma,0)=\chi\exp\bigl(\sum_kz^{2k-1}\mathsf S_{2k}(\sigma)\bigr)$ lies in $\bbA_\sigma\fps{z}$, because $\sigma^{-1}$ and the power sums $\sum_j(\zeta^j-\sigma)^{-r}\in(1-\sigma^5)^{-r}\bbQ[\sigma]$ lie in $\bbA_\sigma$. So $\chi a_0=(\mathsf r(\sigma)/\mathsf r(L))^{1/2}b$, and \eqref{eq:fsp-generic} gives
\begin{equation}
\label{eq:con-shiftD}
(z\Dq+L)\bigl(L\,c\langle\chi\rangle\bigr)=L\,c\langle\sigma\chi\rangle\quad\text{in }\AL\lser{z},\qquad\chi\in\bbQ[\sigma,z].
\end{equation}

\emph{Step 2 ($\PFfive$).} $\PFfive=(z\Dq)^5-1-q\prod_{k=1}^5(5z\Dq+kz)$, with $q$ acting by multiplication on the left. Iterating \eqref{eq:con-shiftD} and using the $\AL[z]$-linearity of $\chi\mapsto c\langle\chi\rangle$,
\[
\PFfive^{L}(L\,c)=L\Bigl(c\langle\sigma^5-1\rangle-c\bigl\langle q\textstyle\prod_{k=1}^5(5\sigma+kz)\bigr\rangle\Bigr).
\]

\emph{Step 3 (the shift).} Put $a:=(\sigma^5-1)a_0$, and let $\mathsf T_z$ be the continuous ring endomorphism $f(t,z)\mapsto f(t+z,z)$ of $\AL\fps{t,z}$. The translate of $a$ by $z$ (\cref{lem:fsp-transl}~(b)) is $b_z=\re^{\mathsf E}\,\mathsf T_z(a)$ with
\[
\mathsf E:=\bigl(\mathsf S_0(L+t+z)-\mathsf S_0(L+t)\bigr)/z=\sum_{n\geq1}\mathsf S_0^{(n)}(L+t)\,z^{n-1}/n!\in\AL\fps{t,z},
\]
which has no constant term, because $\mathsf S_0'(L)=0$. We claim
\begin{equation}
\label{eq:con-bz}
b_z=q\prod_{k=1}^5(5\sigma+kz)\,a_0\qquad\text{in }\AL\fps{t,z},\quad\sigma=L+t .
\end{equation}
The ring $\AL\fps{t,z}$ embeds in $\Bbbk\fps{t,z}$, where $\Bbbk:=\AL\otimes_{\bbQ}\bbQ(\zeta)=\bbQ(\zeta)[L,L^{-1},(L^5-1)^{-1}]$ (in this proof only; $\bbQ(\zeta)$ is free over $\bbQ$, so $\AL\to\Bbbk$ is injective), in which $5L$ and the $\zeta^j-L$ are units; we prove \eqref{eq:con-bz} there. We write $\mathsf T_z$ also for the endomorphism $f(t,z)\mapsto f(t+z,z)$ of $\Bbbk\fps{t,z}$; it is continuous and preserves the ideal $(t,z)$, so it commutes with $\exp$ and $\log$. Let $\rho(\sigma)=\alpha\sigma+\beta$ with $\alpha\in\bbZ$, $\beta\in\bbQ(\zeta)$ and $\rho_0:=\rho(L)\in\Bbbk^\times$, so that $\rho(L+t)=\rho_0+\alpha t$ is a unit of $\Bbbk\fps t$, again written $\rho$, and put
\begin{gather*}
\mathsf l_\rho:=\log(1+\alpha t/\rho_0)\in t\Bbbk\fps t,\qquad\mathsf U_\rho:=\rho\,\mathsf l_\rho-\alpha t\in t^2\Bbbk\fps t,\\
\mathsf m_\rho:=\sum_{k\geq1}\bk_k\,z^{2k-1}\rho^{1-2k},\qquad w:=z/\rho\in z\Bbbk\fps{t,z};
\end{gather*}
these are the Taylor parts at $\sigma=L$ of $\log\rho$ (orders $\geq1$) and of $\rho\log\rho-\rho$ (orders $\geq2$), and the Stirling tail $\mu(\rho/z)$ of \cref{def:fsp-stirling}. Then
\begin{equation}
\label{eq:con-factor}
\frac{\mathsf T_z\mathsf U_\rho-\mathsf U_\rho}z+(\mathsf T_z-1)\bigl(-\tfrac12\mathsf l_\rho+\mathsf m_\rho\bigr)=\alpha\,\mathsf l_\rho+\log\mathrm{Poch}_\alpha(z/\rho)\qquad\text{in }\Bbbk\fps{t,z}.
\end{equation}
Indeed $\mathsf T_z\rho=\rho\,(1+\alpha w)$, so $\mathsf T_z\mathsf l_\rho=\mathsf l_\rho+\log(1+\alpha w)$ and $(\mathsf T_z\mathsf U_\rho-\mathsf U_\rho)/z=\alpha\,\mathsf l_\rho+(w^{-1}+\alpha)\log(1+\alpha w)-\alpha$ (here $w^{-1}\log(1+\alpha w)$ is the power series $\sum_{n\geq1}(-1)^{n+1}\alpha^nw^{n-1}/n$, as in \cref{def:fsp-stirling}), while $\mathsf m_\rho=\sum_k\bk_kw^{2k-1}$ and $\mathsf T_z\mathsf m_\rho=\sum_k\bk_kw^{2k-1}(1+\alpha w)^{1-2k}$. So the left side of \eqref{eq:con-factor} is $\alpha\,\mathsf l_\rho+\log\mathrm{St}_\alpha(w)-\log\mathrm{St}_0(w)$, and $\mathrm{St}_\alpha=\mathrm{Poch}_\alpha\,\mathrm{St}_0$ in $\bbQ\fps w$ by \cref{lem:fsp-stirfe}; the substitution $w=z/\rho$ is a continuous ring homomorphism $\bbQ\fps w\to\Bbbk\fps{t,z}$.

Apply \eqref{eq:con-factor} to $\rho_{\mathrm N}:=5\sigma$ ($\alpha=5$) and to $\rho_j:=\zeta^j-\sigma$ ($\alpha=-1$), $j=0,\dots,4$. Since $[t^{k+1}]\,\mathsf U_\rho=(-1)^{k+1}\alpha^{k+1}\rho_0^{-k}/(k(k+1))$ and $[t^k]\,\mathsf l_\rho=(-1)^{k+1}\alpha^k\rho_0^{-k}/k$ for $k\geq1$, the Taylor data of \cref{def:gam-c} give, coefficient by coefficient, $\mathsf S_0(L+t)-\mathsf S_0(L)=\mathsf U_{\rho_{\mathrm N}}+\sum_j\mathsf U_{\rho_j}$, $\mathsf S_1(L+t)-\mathsf S_1(L)=\frac12\mathsf l_{\rho_{\mathrm N}}-\frac12\sum_j\mathsf l_{\rho_j}$ and $\sum_kz^{2k-1}\mathsf S_{2k}(L+t)=\mathsf m_{\rho_{\mathrm N}}+\sum_j\mathsf m_{\rho_j}$. Summing \eqref{eq:con-factor} over the six factors and adding $(\mathsf T_z-1)\mathsf l_{\rho_{\mathrm N}}=\log(1+5z/\rho_{\mathrm N})$, with $\mathrm{Poch}_5(w)(1+5w)=\prod_{k=1}^5(1+kw)$ and $\mathrm{Poch}_{-1}(w)=(1-w)^{-1}$,
\[
\mathsf E+(\mathsf T_z-1)\log a_0=5\,\mathsf l_{\rho_{\mathrm N}}-\sum_j\mathsf l_{\rho_j}+\sum_{k=1}^5\log\Bigl(1+\frac{kz}{\rho_{\mathrm N}}\Bigr)-\sum_j\log\Bigl(1-\frac z{\rho_j}\Bigr).
\]
Exponentiate, using $\re^{\mathsf l_\rho}=\rho/\rho_0$, $\prod_j(\zeta^j-x)=1-x^5$ and $\prod_j(\zeta^j-L)/(5L)^5=-q$ (\cref{def:gam-c}):
\[
\re^{\mathsf E}\,\frac{\mathsf T_z(a_0)}{a_0}=\frac{\prod_{k=1}^5(5\sigma+kz)}{(5L)^5}\cdot\frac{\prod_j(\zeta^j-L)}{\prod_j(\zeta^j-\sigma-z)}=\frac{q\prod_{k=1}^5(5\sigma+kz)}{(\sigma+z)^5-1}.
\]
Multiplying by $\mathsf T_z(\sigma^5-1)\,a_0=((\sigma+z)^5-1)\,a_0$ gives \eqref{eq:con-bz}. By \cref{lem:fsp-transl}~(b), with $p=z$ and $a\in\AL\fps t\fps z$, $\SPn_L[\mathsf S_0;b_z]=\SPn_L[\mathsf S_0;a]$, that is, $c\langle q\prod_k(5\sigma+kz)\rangle=c\langle\sigma^5-1\rangle$; by Step~2, $\PFfive^{L}(L\,c)=0$.
\end{proof}

% Section 8 (part 5): Theorem P(a), parts (ii)-(iv); the uniformising map; reduction modulo exact forms.

\begin{proof}[Proof of (ii)]
Let $A:=\bbQ\lser{s}$, a field, in which $L=s^{-1}$ and $u=s^5$. By \cref{lem:fsp-univ}, $\jmath(c)$ is the normalised saddle series of the same data computed over $A$. As in the proof of (i), $\mathsf S_0$ stands for the formal phase $\mathsf S_0(L+t)-\mathsf S_0(L)$ and the amplitude is $a_0:=A_{\mathrm{St}}/A_{\mathrm{St}}(L,0)$, both with Taylor data in $\AL$, mapped to $A$ by $\jmath$. The letters $\phi$, $J$, $a$, $a_0$ are local to this proof.

\emph{Step 1 (equivariance).} On $A$, $\Dq=-(1-u)\,\dd/\dd u=-\frac{1-s^5}{5s^4}\frac{\dd}{\dd s}$, and $\Dq(s^{-1})=(1-s^5)/(5s^6)=\jmath((L^6-L)/5)$. A ring homomorphism compatible with the derivations on the generators is compatible on the localisation; so $\jmath\circ\Dq=\Dq\circ\jmath$.

\emph{Step 2 (change of variables).} The map $\phi(X):=L/X$, $\phi(X)-L=L\sum_{m\geq1}(1-X)^m$, has $\phi(1)=L$ and $J:=\phi'(1)=-L$, a unit of $A$. By \cref{lem:fsp-cov} over $A$, $\SPn_L[\mathsf S_0;a]=\SPn_{X=1}[\mathsf S_0\circ\phi;(a\circ\phi)\,\phi'/J]$ for every amplitude $a\in A\fps{\sigma-L}\lser z$.

\emph{Step 3 (the phases agree).} In $A\fps{X-1}$, $\mathsf S_0(\phi(X))-\mathsf S_0(L)=-s^4\bigl(\hat\cW(X;u)-\hat\cW(1;u)\bigr)$. Both sides vanish at $X=1$, and their $X$-derivatives are $-(L/X^2)\mathsf S_0'(L/X)$ and $(L/X^2)\log((1-uX^5)/(1-u))$ (\cref{lem:con-form}~(c)); so it suffices that $\mathsf S_0'(L/X)=-\log((1-uX^5)/(1-u))$. Both sides of this vanish at $X=1$, and both $X$-derivatives are images of $5uX^4/(1-uX^5)$, $u=s^5$ (on the left $\mathsf S_0''(L/X)\cdot(-L/X^2)$, with $\mathsf S_0''=-5/(\sigma(\sigma^5-1))$), under two ring homomorphisms $\bbQ[s^{\pm1},X^{\pm1},(1-s^5X^5)^{-1}]\to A\fps{X-1}$: expansion at $X=1$ over $\bbQ(s)$, and expansion of $(1-s^5X^5)^{-1}=\sum_Ps^{5P}X^{5P}$ in $\bbQ[X,X^{-1}]\lser s$ followed by expansion at $X=1$ coefficientwise in $s$, a ring homomorphism $\bbQ[X,X^{-1}]\lser s\to A\fps{X-1}$ which commutes with $\dd/\dd X$. They agree on the generators, hence everywhere; and a series in $X-1$ is determined by its value at $X=1$ and its derivative.

\emph{Step 4 (the amplitudes agree).} In $A\fps{X-1}\fps{z}$, with $z=s^4\zh$, $(A_{\mathrm{St}}\circ\phi)\,\phi'/(J\,A_{\mathrm{St}}(L,0))=\hat\fa/\hat\fa(1;0)$. The half-density: at $\sigma=L/X$ one has $5\sigma/(1-\sigma^5)=5LuX^4/(uX^5-1)$ and $\phi'/J=X^{-2}$, so $\re^{\mathsf S_1(L/X)-\mathsf S_1(L)}\phi'/J$ is the square root with constant term $1$ of $(1-u)/(1-uX^5)$, which is the half-density of $\hat\fa/\hat\fa(1;0)$. The Stirling parts: $z^{2k-1}(5\sigma)^{1-2k}=5^{1-2k}(uX\zh)^{2k-1}$ and
\[
z^{2k-1}\sum_j(\zeta^j-\sigma)^{1-2k}=-(uX\zh)^{2k-1}\sum_j(1-\zeta^jsX)^{1-2k}=-(uX\zh)^{2k-1}\,5\sum_{P\geq0}\binom{2k-2+5P}{5P}u^PX^{5P},
\]
which gives the exponent of $\hat\fa$ (\cref{lem:con-form}~(b), (c)); these identities hold by the two-homomorphism argument of Step~3, applied to the sums over $j$, which lie in $\bbQ[s,X,(1-s^5X^5)^{-1}]$.

\emph{Step 5 (the rescaling).} Write the amplitude of Step 4 as $a(X;z):=\hat\fa(X;z/s^4)/\hat\fa(1;0)$. By \cref{lem:fsp-homog} with the unit $\mu=s^4$ of $A$, $\SPn_{X=1}[-s^4\hat\cW;a(\cdot;z)](z)=\SPn_{X=1}[-\hat\cW;a(\cdot;s^4\zh)](\zh)$ at $\zh=z/s^4$, and $a(\cdot;s^4\zh)=\hat\fa(\cdot;\zh)/\hat\fa(1;0)$. With Steps 2--4 and \cref{lem:fsp-const} this gives $\jmath(c)(z)=\PhiSP(z/s^4;u)$ in $A\fps{z}$. The hatted data lie in $\bbQ\fps{u}$ (\cref{lem:con-form}~(c)), so $\PhiSP$ is computed in $\bbQ\fps{u}\fps{\zh}$ and, by \cref{lem:fsp-univ} applied to $\bbQ\fps{u}\to A$, equals the same functional over $A$.
\end{proof}

\begin{proof}[Proof of (iii)]
Apply $\jmath$ to (i): by Step 1 of the proof of (ii), $\PFfive^{L}(s^{-1}\jmath(c))=0$ in $\bbQ\lser{s}\fps{z}$. The substitution $z=s^4\zh$ is a ring isomorphism $\bbQ\lser{s}\fps{z}\to\bbQ\lser{s}\fps{\zh}$, since $s^4$ is a unit, and by (ii) it maps $s^{-1}\jmath(c)$ to $s^{-1}\PhiSP$. By \eqref{eq:con-chain}, $\PFfive^{L}(s^{-1}\PhiSP)=s^{-1}\mathrm{HPF}(\PhiSP)$, so $\mathrm{HPF}(\PhiSP)=0$ in $\bbQ\lser{s}\fps{\zh}$. HPF has coefficients in $\bbQ[u,\zh]$, and $\PhiSP$ lies in $\bbQ\fps{u}\fps{\zh}$, which embeds in $\bbQ\lser{s}\fps{\zh}$; so the identity holds in $\bbQ\fps{u}\fps{\zh}$.
\end{proof}

\begin{proof}[Proof of (iv)]
By (iii), $\PhiSP$ solves HPF with $\mathrm{ph}^{\mathrm{SP}}_0=1$ and coefficients in $\bbQ\fps{u}$; by \cref{lem:gau-poly} (only its polynomiality clause is used) every $\mathrm{ph}^{\mathrm{SP}}_n$ is a polynomial. Then $\jmath\bigl(L^{4n}\mathrm{ph}^{\mathrm{SP}}_n(L^{-5})\bigr)=s^{-4n}\mathrm{ph}^{\mathrm{SP}}_n(s^5)=\jmath(c_n)$ by (ii), and $\jmath$ is injective.
\end{proof}

Theorem P(b) identifies $\PhiSP$ with $\Phican$. Its proof uniformises the hatted phase by the model phase $\Wv$, reduces the $1$-form modulo exact forms to the five model forms, and compares two functionals on the reduced form, the stationary phase at the critical point and the residue at the pole.

\begin{lem}[The uniformising map]\genusstatus{AI-proposed, unchecked}
\label{lem:con-cfu}
There are unique $Y_k\in\bbQ[T]$ ($k\geq1$) and $V\in1+u\bbQ\fps{u}$ such that $Y:=X\bigl(1+\sum_{k\geq1}u^kY_k(X^5)\bigr)$ satisfies
\[
\hat\cW(X;u)=V/Y+Y^4/4=W_V(Y)\qquad\text{in }\bbQ[X,X^{-1}]\fps{u}.
\]
Moreover $\deg Y_k\leq k$ and $Y_1=(1+T)/18$; $V=\tau^{-5/4}$; the point $X=1$ goes to the critical point $Y_c:=Y(1)=V^{1/5}=\tau^{-1/4}$ of $W_V$, and $Y'(1)\in1+u\bbQ\fps{u}$; the inverse map has the same form, $X=Y\bigl(1+\sum_ku^kX_k(Y^5)\bigr)$ with $X_k\in\bbQ[T]$ and $\deg X_k\leq k$; and $Y(\zeta X)=\zeta Y(X)$.
\end{lem}

\begin{proof}
Write $Y=X\,y(T)$, $T=X^5$, $y=1+\sum_{k\geq1}u^kY_k(T)$. Then $V/Y+Y^4/4=X^{-1}\bigl(V/y+Ty^4/4\bigr)$, and by \cref{lem:con-form}~(c) the equation becomes $V/y+Ty^4/4=\sum_{m\geq1}\frac{u^{m-1}}m\bigl(1+\frac{T^m}{5m-1}\bigr)$. At $u^0$ it holds with $V=1$, $y=1$. At $u^k$, $k\geq1$, it reads $(T-1)Y_k+V_k=R_k(T)$, where $V_k:=[u^k]V$ and $R_k$ is a polynomial built from $Y_1,\dots,Y_{k-1}$, $V_1,\dots,V_{k-1}$ and the term $\frac1{k+1}(1+\frac{T^{k+1}}{5k+4})$. It has a solution with $Y_k$ a polynomial if and only if $V_k=R_k(1)$, and then $Y_k=(R_k(T)-R_k(1))/(T-1)$ is unique. By induction $\deg Y_j\leq j$ for $j<k$; the $u^k$-parts of $V/y$ and of $Ty^4/4$ other than $V_k-Y_k$ and $TY_k$ then have degree at most $k$ and $k+1$, so $\deg R_k\leq k+1$ and $\deg Y_k\leq k$. At $k=1$: $R_1=\frac12+\frac{T^2}{18}$, $V_1=\frac59$, $Y_1=(1+T)/18$. Since $\hat\cW'(1)=0$ and $Y'(1)\in1+u\bbQ\fps{u}$ is a unit, $W_V'(Y(1))=0$, i.e. $Y(1)^5=V$ with $Y(1)\in1+u\bbQ\fps{u}$, so $Y(1)=V^{1/5}$; then $W_V(Y_c)=\frac54V^{4/5}$, and $\hat\cW(1)=\frac54\tau^{-1}$ gives $V=\tau^{-5/4}$. In the variables $x=sX$ and $sY$ the map reads $sY=x\bigl(1+\sum c_{k,d}u^{k-d}x^{5d}\bigr)$ with $d\leq k$, a change of variables in $\bbQ\fps{u,x}$ of the form $x\,(1+g(u,x^5))$ with $g(0,0)=0$. Its compositional inverse exists in $\bbQ\fps{u,y}$ and commutes with $y\mapsto\zeta y$, so it has the form $y\,(1+h(u,y^5))$ with $h(0,0)=0$; a monomial $u^ay^{5d}$ of $h$ is $u^{a+d}Y^{5d}$ at $y=sY$, and returning to $X$, $Y$ gives the stated form of the inverse with $\deg X_k\leq k$. Finally $T=X^5$ is invariant under $X\mapsto\zeta X$.
\end{proof}

\begin{lem}[Reduction modulo exact forms]\genusstatus{AI-proposed, unchecked}
\label{lem:con-red}
Let $\fa_Y(Y;\zh):=\hat\fa(X(Y);\zh)\,X'(Y)$, the hatted amplitude in the variable $Y$ of \cref{lem:con-cfu}; it lies in $\bbQ[Y,\zh]\fps{u}$. For $n\geq4$,
\[
\zh\,\dd\bigl(Y^{n-3}\re^{-W_V/\zh}\bigr)=\bigl((n-3)\zh Y^{n-4}+VY^{n-5}-Y^n\bigr)\re^{-W_V/\zh}\,\dd Y,
\]
and there are $\pi_m\in\bbQ[\zh]\fps{u}$ ($m=-1,\dots,3$) and $\mathcal E\in Y\bbQ[Y,\zh]\fps{u}$ with
\[
\fa_Y\,\re^{-W_V/\zh}\,\dd Y=\sum_{m=-1}^{3}\pi_m\,Y^m\re^{-W_V/\zh}\,\dd Y+\dd\bigl(\mathcal E\,\re^{-W_V/\zh}\bigr).
\]
With $\deg Y=1$ and $\deg\zh=-1$, $\pi_m$ contains only powers $\zh^k$ with $k\equiv m\bmod5$; for $m=\nu_b$ this is $k\equiv-b\bmod5$.
\end{lem}

\begin{proof}
The displayed identity is $\frac{\dd}{\dd Y}(Y^{n-3}\re^{-W_V/\zh})=\bigl((n-3)Y^{n-4}-Y^{n-3}W_V'(Y)/\zh\bigr)\re^{-W_V/\zh}$ with $W_V'=Y^3-VY^{-2}$; so $Y^n\equiv(n-3)\zh Y^{n-4}+VY^{n-5}$ modulo exact forms. At each $u$-order, $\fa_Y$ is a polynomial in $Y$ and $\zh$, because $\hat\fa\in\bbQ[X,\zh]\fps{u}$ and the inverse map of \cref{lem:con-cfu} is polynomial at each $u$-order. Each reduction step lowers the $Y$-degree by at least $4$, so finitely many steps bring a polynomial to the exponents $-1,\dots,3$, the lowest exponent produced being $4-5=-1$; the powers of $V$ bring power series in $u$, and each $u$-order receives finitely many steps. The primitives are $Y^{n-3}$ with $n\geq4$, whence $\mathcal E\in Y\bbQ[Y,\zh]\fps{u}$. Grading: $W_V(\zeta Y)=\zeta^{-1}W_V(Y)$; $\hat\fa(\zeta X;\zh)=\hat\fa(X;\zeta\zh)$, because $\fa$ depends on $x$ only through $xz$ and $x^5$; and the uniformising map commutes with $X\mapsto\zeta X$. So every monomial $Y^a\zh^b\,\dd Y$ of $\fa_Y\,\dd Y$ has $a-b+1\equiv1\bmod5$, and the relation preserves the degree modulo $5$. Hence every monomial $\zh^kY^m\,\dd Y$ of the result has $m-k\equiv0$; and $\nu_b\equiv-b\bmod5$ (\cref{def:set-model}). Uniqueness of the $\pi_m$ is not needed: \cref{lem:con-id} below determines them.
\end{proof}

% Section 8 (part 6): the model identity; signs, weights and identification; the relabelling.

\begin{lem}[The model identity]\genusstatus{AI-proposed, unchecked}
\label{lem:con-S2}
Let $W_1(\tilde p):=\tilde p^4/4+1/\tilde p$. For every $b\in\bbZ/5$, $\SPn_{\tilde p=1}[-W_1;\tilde p^{\nu_b}](t)=\varepsilon_b\Fs_b(t)$; equivalently, the full stationary phase $\mathrm{SP}_{\tilde p=1}[\tilde p^{\nu_b}\re^{-W_1(\tilde p)/t}\,\dd\tilde p]$ divided by $(2\pi t/5)^{1/2}\re^{-5/(4t)}$ is $\varepsilon_b\Fs_b(t)$.
\end{lem}

\begin{proof}
(The letters $J_m$, $G_b$, $U$ are local.) Let $A:=\bbQ\fps{v-1}$ and $\varsigma:=v^{1/5}\in A$. The phase $-\Wv(p)$ has the critical point $p=\varsigma$ (since $\Wv'(p)=p^3-v/p^2$) with $\Wv''(\varsigma)=5\varsigma^2$, a unit of $A$. Let $J_m:=\mathrm{SP}_{p=\varsigma}[p^m\re^{-\Wv(p)/z}\,\dd p]$, the full stationary phase over $A$. By \cref{lem:fsp-dq}, with the moving critical point $\varsigma$, $z\,\dd J_m/\dd v=-J_{m-1}$; by \cref{lem:fsp-ibp} applied to $\dd(\re^{-\Wv/z})$, $J_3=vJ_{-2}$. Hence $(J_0,-J_{-1},-J_3,J_2,J_1)=(\varepsilon_bJ_{\nu_b})_b$ solves the model system $z\,\dd Y_b/\dd v=b^{\mathrm m}_bY_{b+1}$, $b^{\mathrm m}=(1,-1/v,1,-1,-1)$ (\cref{prop:set-modrows}); the row $b=1$ uses $J_{-2}=J_3/v$, the other rows the derivative identity alone. Next, \cref{lem:fsp-cov} over $A$ with the linear map $p=\varsigma\tilde p$, for which $\Wv(\varsigma\tilde p)=\varsigma^4W_1(\tilde p)$, followed by \cref{lem:fsp-homog} with the unit $\mu=\varsigma^4$, gives
\[
J_{\nu_b}=(2\pi z/5)^{1/2}\re^{U/z}\varsigma^{\nu_b-1}G_b(z/\varsigma^4),\qquad U=-\tfrac54\varsigma^4,
\]
where $G_b(t):=\SPn_{\tilde p=1}[-W_1;\tilde p^{\nu_b}](t)$ does not depend on $v$. With $\dd U/\dd v=-1/\varsigma$ and $\dd\varsigma/\dd v=1/(5\varsigma^4)$ the chain rule gives $z\frac{\dd}{\dd v}\bigl[\re^{U/z}\varsigma^{k}F(z/\varsigma^4)\bigr]=\re^{U/z}\varsigma^{k-1}\mathcal T_kF$ with $\mathcal T_k=-1+\frac t5(k-4t\frac{\dd}{\dd t})$, so the system becomes the colour-$0$ T-system of \cref{prop:set-modrows} for $F_b:=\varepsilon_bG_b$. Since $G_0(0)=1$, the uniqueness half of \cref{prop:set-modrows} (pivot $4n$) gives $F_b=\Fs_b$, i.e. $G_b=\varepsilon_b\Fs_b$.
\end{proof}

\begin{lem}[Signs, weights, identification]\genusstatus{AI-proposed, unchecked}
\label{lem:con-id}
Work in a conifold frame of \cref{def:set-frames}, and let $Y_c=\tau^{-1/4}$ and the $\pi_m$ be as in \cref{lem:con-cfu,lem:con-red}.
\begin{enumerate}
\item[(a)] $\mathrm{SP}_{Y_c}[Y^{\nu_b}\re^{-W_V/\zh}\,\dd Y]=N_{\zh}\,\varepsilon_b\,\tau^{(1-\nu_b)/4}\Fs_b(\tau\zh)$, where $N_{\zh}:=(2\pi\zh/5)^{1/2}\re^{-W_V(Y_c)/\zh}$ and $W_V(Y_c)=\hat\cW(1)$.
\item[(b)] $\PhiSP=\sum_b\varepsilon_b\,\pi_{\nu_b}\,\tau^{(1-\nu_b)/4}\Fs_b(\tau\zh)$.
\item[(c)] Let $P$ be a graded regular gauge with coefficients in $\bbQ\fps u$ from the pulled-back model to the quintic system (\cref{def:gau-onepole,prop:gau-graded}), and let $\Phi_P$ be $\varpi_0$ times the colour-$0$ entry of the hatted row $0$ of $P\,\Qr^{\mathrm m}(v(u))$; for $P=\Pc$ this is $\Phican$ (\cref{thm:gau-R0}). Then
\[
\Phi_P/\varpi_0=\sum_b\tau^{(1-\nu_b)/4}\,G_b\,\Fs_b(\tau\zh),\qquad G_b:=s^{\nu_b}P_{0b}(u,s^4\zh),
\]
where $G_b\in\bbQ\fps{u}\fps{\zh}$ contains only powers $\zh^k$ with $k\equiv-b\bmod5$, and $[\zh^k]G_b$ is divisible by $u^{(\nu_b+4k)/5}$.
\item[(d)] Let $\Delta_b\in\bbQ\fps{u}\fps{\zh}$, $b\in\bbZ/5$, be series that contain only powers $\zh^k$ with $k\equiv-b\bmod5$. If $\sum_b\tau^{(1-\nu_b)/4}\Delta_b\Fs_b(\tau\zh)=0$, then every $\Delta_b$ vanishes.
\end{enumerate}
Hence, if $P$ is a graded regular gauge with $\Phi_P=\PhiSP$, then $G_b=\varepsilon_b\pi_{\nu_b}/\varpi_0$ for every $b$.
\end{lem}

\begin{proof}
(a) Put $Y=Y_c\tilde p$. Since $V=Y_c^5$, $W_V(Y_c\tilde p)=Y_c^4W_1(\tilde p)$, and $\zh/Y_c^4=\tau\zh$. By \cref{lem:fsp-cov} (a linear map over $\bbQ\fps{u}$) and \cref{lem:fsp-homog} with $\mu=Y_c^4=\tau^{-1}$, the left side is $Y_c^{\nu_b+1}\,\mathrm{SP}_{\tilde p=1}[\tilde p^{\nu_b}\re^{-W_1/t}\,\dd\tilde p]$ at $t=\tau\zh$, which by \cref{lem:con-S2} is $Y_c^{\nu_b+1}(2\pi\tau\zh/5)^{1/2}\re^{-5/(4\tau\zh)}\varepsilon_b\Fs_b(\tau\zh)$. Here $Y_c^{\nu_b+1}\tau^{1/2}=\tau^{(1-\nu_b)/4}$ and $\frac5{4\tau}=\frac54Y_c^4=W_V(Y_c)=\hat\cW(1)$ (\cref{lem:con-cfu}).

(b) By \cref{lem:fsp-cfu} the full stationary phase at $X=1$ of $\re^{-\hat\cW/\zh}\hat\fa\,\dd X$ equals that at $Y_c$ of its image $\re^{-W_V/\zh}\fa_Y\,\dd Y$. By \cref{lem:con-red} and \cref{lem:fsp-ibp}, applied once to the exact form $\dd(\mathcal E\re^{-W_V/\zh})$, whose primitive $\mathcal E$ is one amplitude over $\bbQ\fps{u}$, and by the $\bbQ\fps{u}\fps{\zh}$-linearity of the functional (\cref{def:fsp-wick}), this is $\sum_m\pi_m\,\mathrm{SP}_{Y_c}[Y^m\re^{-W_V/\zh}\,\dd Y]$. The exponents $m=-1,\dots,3$ are the $\nu_b$; insert (a) and divide by $N_{\zh}=(2\pi\zh/5)^{1/2}\re^{-\hat\cW(1)/\zh}$ (\cref{def:con-PhiSP}). The square roots agree: \cref{lem:fsp-cfu} takes $\hat\cW''(1)^{1/2}=Y'(1)\,W_V''(Y_c)^{1/2}$, (a) takes $W_V''(Y_c)^{1/2}=\sqrt5\,Y_c$, and \cref{def:con-PhiSP} takes $\hat\cW''(1)^{1/2}=\sqrt5\,(1-u)^{-1/2}$. Differentiating $\hat\cW=W_V\circ Y$ twice at $X=1$, where $W_V'(Y_c)=0$, gives $5/(1-u)=5Y_c^2\,Y'(1)^2$; as $Y_c$, $Y'(1)$ and $(1-u)^{-1/2}$ lie in $1+u\bbQ\fps u$, where a square root with constant term $1$ is unique, $Y'(1)\,Y_c=(1-u)^{-1/2}$.

(c) The colour-$0$ model column is $\Qr^{\mathrm m}_b=\varsigma^{\nu_b-1}\Fs_b(z/\varsigma^4)$ after the common exponential is stripped (\cref{prop:set-modrows,thm:gau-R0}), and $\varsigma=s\tau^{-1/4}$, $z/\varsigma^4=\tau\zh$ (\cref{def:set-tau}). Row $0$ of $P\Qr^{\mathrm m}$, hatted by $s^{1-\nu_0}=s$, is $\sum_bs\,P_{0b}\,s^{\nu_b-1}\tau^{(1-\nu_b)/4}\Fs_b(\tau\zh)$, which is the stated formula. For graded $P=\sum_kz^k\sum_{b'}p_{k,b'}E_{b'+k,b'}$ the entry $P_{0b}$ contains only $z^k$ with $k\equiv-b$, and $s^{\nu_b}z^k=u^{(\nu_b+4k)/5}\zh^k$ with $\nu_b+4k\equiv-b-k\equiv0\bmod5$ and $\nu_b+4k\geq0$.

(d) If not all $\Delta_b$ vanish, let $N$ be the least $\zh$-degree occurring in them. Only $b\equiv-N$ can contribute to $\zh^N$, through $\tau^{(1-\nu_b)/4}\varepsilon_b[\zh^N]\Delta_b$, because $\Fs_b(0)=\varepsilon_b$; as $\tau$ is a unit and $\varepsilon_b=\pm1$, this coefficient vanishes, a contradiction.

Finally, if $\Phi_P=\PhiSP$, subtract (b) from $\varpi_0$ times (c): the family $\Delta_b:=G_b-\varepsilon_b\pi_{\nu_b}/\varpi_0$ is graded as in (d), by (c) and \cref{lem:con-red} ($\nu_b\equiv-b$), so it vanishes.
\end{proof}

\begin{lem}[The relabelling]\genusstatus{AI-proposed, unchecked}
\label{lem:con-relab}
Let $\hat\cV:=\Res_X[\re^{-\hat\cW/\zh}\hat\fa\,\dd X]$, computed coefficientwise in $u$ in the residue ring $\cR_{\mathrm{res}}$ of \cref{lem:fsp-res}; at each $u$-order the integrand is $\re^{-(1/X+X^4/4)/\zh}$ times a Laurent polynomial in $X$ over $\bbQ[\zh,\zh^{-1}]$, an element of $\cR_{\mathrm{res}}$. Then $\cV=s\,\hat\cV$ coefficientwise: $[u^N\zh^j]\hat\cV=[u^{N+(1-4j)/5}z^j]\cV$ for all $N$ and $j$.
\end{lem}

\begin{proof}
\emph{Labels.} Monomials are $u^Ez^\beta x^\alpha$ in the flat variables and $u^N\zh^jX^\alpha$ in the hatted ones; the substitution $x=sX$, $z=s^4\zh$, $u=s^5$ acts on exponents by $\iota(E,\beta,\alpha)=(E+(4\beta+\alpha)/5,\beta,\alpha)$, which is injective. With $\dd x=s\,\dd X$ the claim is the coefficientwise statement.
\emph{The family.} The logarithm of $\re^{-\cG/z}\fa$ is a sum of atoms $-x^{5m-1}/(m(5m-1)z)$, $x^{5n}/(2n)$ and $\bk_\kappa\,c_{\kappa,P}\,(xz)^{2\kappa-1}x^{5P}$ ($\kappa\geq1$ a local index), with $c_{\kappa,P}:=5^{1-2\kappa}\delta_{P,0}-5\binom{2\kappa-2+5P}{5P}$ (\cref{lem:con-form}~(b)). A member of the family is a triple ($k$, $r$, a finite multiset of atoms); its term is $[u^r](\Xi/u)^k\,\frac{(-1)^k}{k!}\,u^{k+r}(xz)^{-k}$ times the product of the atoms divided by the factorials of the multiplicities. By \cref{def:con-PhiSP}, $\cV$ is the sum of the $x^{-1}$-parts of these terms.
\emph{Flat fibres.} A fixed $(u^E,z^\beta,x^{-1})$ receives finitely many members: $k\leq E$, and every atom has positive $x$-degree, so the multiset is bounded by the $x$-degree $k-1$.
\emph{Hatted fibres.} Every atom $x^az^b$ has hatted $u$-order $(a+4b)/5\geq0$, the only atom of hatted order $0$ is $-x^4/(4z)\mapsto-X^4/(4\zh)$, and $u^{k+r}(xz)^{-k}\mapsto u^r(X\zh)^{-k}$. Fix $(u^N,\zh^j,X^{-1})$. Then $r\leq N$, at most $N$ atoms have positive hatted order, from a finite list, and the $X$- and $\zh$-balances $-k+4l+a'=-1$, $-k-l+b'=j$, where $l$ is the multiplicity of $X^4/(4\zh)$ and $a'$, $b'$ are the degrees of the other atoms, determine $k$ and $l$ (the determinant is $5$). So this fibre is finite too.
\emph{$\hat\cV$ is the hatted sum.} In $\cR_{\mathrm{res}}\fps{u}$ the product is the convergent sum of the products of coefficients, and $\re^{-(\Xi/u)/(X\zh)}=\re^{-1/(X\zh)}\re^{-(\Xi/u-1)/(X\zh)}$ holds coefficient by coefficient: at $(u^N,X^{-n})$ both sides equal $(-1/\zh)^n[u^N](\Xi/u)^n/n!$, a finite expression. So $\hat\cV$ is the sum of the family in the hatted order.
\emph{Conclusion.} $\cV$ and $s\hat\cV$ are sums of one family over the fibres of two injective labellings related by $\iota$, with finite fibres in both orders; so they agree coefficient by coefficient.
\end{proof}

\input{sections/s08_connection_7}



